\documentclass[10pt, aspectratio=169]{beamer}
\input{../../_en_preambulo_beamer}
% Comandos y macros estadísticas compartidas para las presentaciones Beamer en inglés (Metropolis)

% Conjuntos numéricos básicos
\newcommand{\R}{\mathbb{R}}
\newcommand{\N}{\mathbb{N}}
\newcommand{\Q}{\mathbb{Q}}
\newcommand{\Z}{\mathbb{Z}}
\newcommand{\set}[1]{\left\{ #1 \right\}}
\newcommand{\abs}[1]{\left|#1\right|}
\newcommand{\norm}[1]{\left\| #1 \right\|}

% Comandos estadísticos y probabilísticos del curso
\newcommand{\Var}{\operatorname{Var}}
\newcommand{\cov}{\operatorname{Cov}}
\newcommand{\corr}{\rho}
\newcommand{\comb}[2]{\begin{pmatrix} #1 \\ #2 \end{pmatrix}}
\newcommand{\s}{\sigma}
\newcommand{\particion}{\mathfrak{P}}
\newcommand{\card}[1]{\#\left( #1 \right)}
\newcommand{\seq}[3]{\set{#1}_{#2}^{#3}}
\newcommand{\E}{\mathbb{E}}
\newcommand{\Prob}{\mathbb{P}}

% Entornos pedagógicos y cajas en inglés académico natural (soportando tanto nombres en inglés como en español)
\newenvironment{discussion}{%
  \begin{alertblock}{Discussion Question}%
}{%
  \end{alertblock}%
}
\newenvironment{discusion}{%
  \begin{alertblock}{Discussion Question}%
}{%
  \end{alertblock}%
}

\newenvironment{reflection}{%
  \begin{alertblock}{Classroom Reflection}%
}{%
  \end{alertblock}%
}
\newenvironment{reflexion}{%
  \begin{alertblock}{Classroom Reflection}%
}{%
  \end{alertblock}%
}

\newenvironment{paradox}{%
  \begin{alertblock}{Exploratory Paradox}%
}{%
  \end{alertblock}%
}
\newenvironment{paradoja}{%
  \begin{alertblock}{Exploratory Paradox}%
}{%
  \end{alertblock}%
}

\newenvironment{motivation}{%
  \begin{exampleblock}{Motivation: Data Science in Practice}%
}{%
  \end{exampleblock}%
}
\newenvironment{motivacion}{%
  \begin{exampleblock}{Motivation: Data Science in Practice}%
}{%
  \end{exampleblock}%
}

\newenvironment{quickchallenge}{%
  \begin{block}{Quick Team Challenge}%
}{%
  \end{block}%
}
\newenvironment{retoclase}{%
  \begin{block}{Quick Team Challenge}%
}{%
  \end{block}%
}


\title[Multiple Regression]{Section 09.07: Multiple Linear Regression}
\subtitle{The Normal Equation in Matrix Notation and Ridge/Lasso Regularization}
\author[J. Castillo Colmenares]{Juliho Castillo Colmenares}
\institute{Tecnológico de Monterrey}
\date{\vspace{-1.5cm}}

\begin{document}

% Slide 1: Title
\begin{frame}[plain]
  \titlepage
\end{frame}

% Slide 2: Roadmap
\begin{frame}{Roadmap --- Chapter 09: Linear and Multiple Regressions}
  \scriptsize
  Where are we in the modular development of this chapter?
  \begin{itemize}\itemsep=0.02em
    \item \textbf{09.01} Correlation as the Premise of Regression
    \item \textbf{09.02} Introduction to Linear Regression
    \item \textbf{09.03} The Mathematics of Regression: OLS Derivation and $R^2$
    \item \textbf{09.04} Linear Regression on Simulated Data
    \item \textbf{09.05} Optimal Coefficients, $t$/$F$ Tests, and RSE
    \item \textbf{09.06} Implementation in Python with \texttt{statsmodels}
    \item \textbf{09.07} Multiple Regression, the Normal Equation, and Ridge/Lasso \emph{(Today)}
    \item \textbf{09.08} Model Validation and $k$-fold Cross-Validation
    \item \textbf{09.09} Residual Diagnostics and Multicollinearity (VIF)
    \item \textbf{09.10} Regression with \texttt{scikit-learn}
    \item \textbf{09.11} Categorical Variables and Dummy Variables
    \item \textbf{09.12} Nonlinear Transformations and Polynomial Regression
  \end{itemize}
  \pause
  \vspace{-0.05cm}
  \begin{block}{Session Objective}
    Generalize simple regression to $p$ simultaneous predictors via matrix notation: derive the Normal Equation $\hat{\boldsymbol\beta}=(\mathbf{X}^T\mathbf{X})^{-1}\mathbf{X}^T\mathbf{Y}$, and present Ridge and Lasso regularization as a remedy when the design matrix is close to singular due to multicollinearity.
  \end{block}
\end{frame}

% Slide 3: Motivation
\begin{frame}{From One Predictor to $p$ Simultaneous Predictors}
  \footnotesize
  \begin{motivacion}
    In practice a single variable rarely explains a phenomenon: a vehicle's fuel efficiency depends on its weight \emph{and} its engine size; salary depends on experience \emph{and} education level. Writing and solving the normal equations with scalar summations for each $\beta_j$ becomes unmanageable as $p$ grows.
  \end{motivacion}

  \vspace{0.15cm}
  \pause
  \begin{itemize}\itemsep=0.1cm
    \item Matrix notation solves the problem compactly and exactly for any $p$. \pause
    \item The model, the objective function, and the solution have the \textbf{same structure} as in simple regression, only expressed with vectors and matrices.
  \end{itemize}
\end{frame}

% Slide 4: Matrix model
\begin{frame}{The Multiple Linear Model in Matrix Form}
  \footnotesize
  \begin{block}{Population Model}
    $$\mathbf{Y}=\mathbf{X}\boldsymbol\beta+\boldsymbol\varepsilon, \qquad E[\boldsymbol\varepsilon]=\mathbf{0},\ \text{Var}(\boldsymbol\varepsilon)=\sigma^2\mathbf{I}_n$$
  \end{block}
  \pause

  \vspace{0.1cm}
  \begin{alertblock}{Dimensions}
    $\mathbf{Y}$ is $n\times1$; $\mathbf{X}$ is $n\times(p+1)$ (first column of ones); $\boldsymbol\beta$ is $(p+1)\times1$. The objective function is the quadratic form $S(\boldsymbol\beta)=(\mathbf{Y}-\mathbf{X}\boldsymbol\beta)^T(\mathbf{Y}-\mathbf{X}\boldsymbol\beta)$.
  \end{alertblock}
\end{frame}

% Slide 5: Normal equation
\begin{frame}{Gauss's Normal Equation}
  \footnotesize
  \begin{block}{Closed-Form Solution}
    $$\mathbf{X}^T\mathbf{X}\hat{\boldsymbol\beta}=\mathbf{X}^T\mathbf{Y} \implies \hat{\boldsymbol\beta}=(\mathbf{X}^T\mathbf{X})^{-1}\mathbf{X}^T\mathbf{Y}$$
  \end{block}
  \pause

  \vspace{0.1cm}
  \begin{alertblock}{Hat Matrix}
    $\mathbf{H}=\mathbf{X}(\mathbf{X}^T\mathbf{X})^{-1}\mathbf{X}^T$ projects $\mathbf{Y}$ onto $\hat{\mathbf{Y}}=\mathbf{H}\mathbf{Y}$. It is symmetric ($\mathbf{H}^T=\mathbf{H}$) and idempotent ($\mathbf{H}^2=\mathbf{H}$): an exact orthogonal projection onto the column space of $\mathbf{X}$.
  \end{alertblock}
\end{frame}

% Slide 6: Regularization motivation
\begin{frame}{When $\mathbf{X}^T\mathbf{X}$ Is Nearly Singular}
  \footnotesize
  \begin{block}{The Multicollinearity Problem}
    If two or more predictors are highly correlated, $\mathbf{X}^T\mathbf{X}$ approaches singularity: its inversion becomes numerically unstable and the variances of $\hat{\boldsymbol\beta}$ explode (formally diagnosed via \texttt{VIF} in Section 09.09).
  \end{block}
  \pause

  \vspace{0.1cm}
  \begin{alertblock}{The Idea Behind Regularization}
    Penalizing the magnitude of the coefficients in the objective function stabilizes the solution in exchange for introducing controlled bias --- the \textbf{bias-variance tradeoff}.
  \end{alertblock}
\end{frame}

% Slide 7: Ridge and Lasso
\begin{frame}{Ridge ($L_2$) and Lasso ($L_1$) Regression}
  \footnotesize
  \begin{block}{Ridge}
    $$S_{\text{Ridge}}=\|\mathbf{Y}-\mathbf{X}\boldsymbol\beta\|^2+\lambda\|\boldsymbol\beta\|_2^2 \implies \hat{\boldsymbol\beta}_{\text{Ridge}}=(\mathbf{X}^T\mathbf{X}+\lambda\mathbf{I})^{-1}\mathbf{X}^T\mathbf{Y}$$
  \end{block}
  \pause

  \vspace{0.1cm}
  \begin{alertblock}{Lasso}
    $$S_{\text{Lasso}}=\|\mathbf{Y}-\mathbf{X}\boldsymbol\beta\|^2+\lambda\|\boldsymbol\beta\|_1 \quad \text{(no closed form; induces \emph{sparsity})}$$
  \end{alertblock}
\end{frame}

% Slide 8: Python Lab (Block 1)
\begin{frame}[fragile]{Python Lab: Normal Equation and Hat Matrix}
  \scriptsize
  Matrix solution $\hat{\boldsymbol\beta}$ verified against \texttt{np.linalg.lstsq}, and numerical verification of the symmetry/idempotence of $\mathbf{H}$ (\texttt{09.02\_multiple\_linear\_regression.py}):
  {\lstset{basicstyle=\fontsize{3.4pt}{4.1pt}\selectfont\ttfamily,aboveskip=0.05em,belowskip=0.05em}
  \lstinputlisting[firstline=18, lastline=40]{../../code/09_regresiones/09.07_multiple_linear_regression.py}}
\end{frame}

% Slide 9: Python Lab (Block 2)
\begin{frame}[fragile]{Python Lab: Ridge Regression}
  \scriptsize
  Closed-form Ridge solution for increasing $\lambda$, observing coefficient shrinkage:
  {\lstset{basicstyle=\fontsize{3.8pt}{4.5pt}\selectfont\ttfamily,aboveskip=0.05em,belowskip=0.05em}
  \lstinputlisting[firstline=43, lastline=51]{../../code/09_regresiones/09.07_multiple_linear_regression.py}}
\end{frame}

% Slide 10: Python Lab (Block 3)
\begin{frame}[fragile]{Python Lab: Lasso via Coordinate Descent}
  \scriptsize
  Manual implementation of the soft-thresholding operator for Lasso:
  {\lstset{basicstyle=\fontsize{3.2pt}{3.9pt}\selectfont\ttfamily,aboveskip=0.05em,belowskip=0.05em}
  \lstinputlisting[firstline=54, lastline=76]{../../code/09_regresiones/09.07_multiple_linear_regression.py}}
\end{frame}

% Slide 11: Terminal Output
\begin{frame}[fragile]{Python Lab: Terminal Output}
  \scriptsize
  Standard output produced when running the lab script:
  \vspace{0.02cm}
  \begin{block}{Standard Console Output}
    \fontsize{3.6pt}{4.3pt}\selectfont
    \begin{verbatim}
=== Block 1: Normal Equation and Hat Matrix ===
beta_hat = [3.1001, 2.1533, -0.5180] (matches np.linalg.lstsq)
Symmetry error=3.26e-16, Idempotence error=4.09e-16

=== Block 2: Ridge Shrinkage ===
lambda=0: [3.10, 2.15, -0.52]  lambda=100: [0.09, 1.41, 0.10]

=== Block 3: Lasso Sparsity ===
lambda=0:    zero coefficients=0/2
lambda=200:  zero coefficients=1/2 (x2 dropped)
lambda=1000: zero coefficients=2/2 (only intercept survives)
    \end{verbatim}
  \end{block}
\end{frame}


% Slide 16: Closing
\begin{frame}{Section Closing: Multiple Linear Regression}
  \begin{reflexion}
    We have generalized simple regression to $p$ predictors via matrix notation: the Normal Equation $\hat{\boldsymbol\beta}=(\mathbf{X}^T\mathbf{X})^{-1}\mathbf{X}^T\mathbf{Y}$ solves the entire model in a single step, and Ridge/Lasso stabilize that solution when multicollinearity threatens to make the matrix nearly singular.
  \end{reflexion}

  \vspace{0.2cm}
  \pause
  \begin{block}{Key Takeaways}
    \begin{itemize}\itemsep=0.08cm
      \item \textbf{Normal Equation:} $\mathbf{X}^T\mathbf{X}\hat{\boldsymbol\beta}=\mathbf{X}^T\mathbf{Y}$, an exact generalization of Section 09.03's scalar formula.
      \item \textbf{Ridge:} $L_2$ penalty, closed-form solution, shrinks coefficients without eliminating them.
      \item \textbf{Lasso:} $L_1$ penalty, no closed form, induces \emph{sparsity} (automatic variable selection).
    \end{itemize}
  \end{block}
\end{frame}

% Slide 17: Modular Perspective
\begin{frame}{Modular Perspective --- The Next Challenge}
  \scriptsize
  \begin{retoclase}
    With the multiple model fitted, the natural question is: does it generalize to new data, or does it merely memorize the training sample? Section 09.08 introduces \textbf{model validation}: the train/test split and $k$-fold cross-validation, the standard tools for estimating out-of-sample error before auditing the model's internal assumptions (Section 09.09) or moving on to its professional implementation (Sections 09.10-09.12).
  \end{retoclase}

  \vspace{0.08cm}
  \pause
  \begin{alertblock}{Recommended Preparation}
    \begin{itemize}\itemsep=0.06cm
      \item Reflect on the difference between fitting error (in-sample) and prediction error (out-of-sample).
      \item Review why evaluating a model on the same data it was trained on produces an optimistic, biased estimate of its real performance.
    \end{itemize}
  \end{alertblock}
\end{frame}

\end{document}
